% !TeX root = ../../Book5.tex
% 纲要标题：### 7.2 Frobenius 特征映射
% 纲要位置：工作记录/计划纲要.md，第 4280 行。
% * 对称群表示环与对称函数；明确整系数 Schur 函数子环及 Frobenius 分次环同构
% * 诱导乘积与对称函数乘法
% * 特征标计算的组合化
% * 对合交换完全与初等对称函数；对应符号扭曲与共轭分拆，供竖条 Pieri 规则使用
\section{Frobenius 特征映射}
\label{b5:sec:0702}

幂和函数适合记录轮换长度：一个 \(r\) 轮换对应一个 \(p_r\)，互不相交的轮换对应乘积。然而同一循环型的共轭类大小不同，直接把特征标值作为系数会破坏内积。中心化子阶 \(z_\mu\) 正好给出所需的归一化。

\subsection{对称群表示环与对称函数}
\label{b5:subsec:0702:01}

\begin{definition}\label{b5:def:frobenius-characteristic}
	设 \(n\) 为非负整数，\(f\) 为 \(S_n\) 上的复类函数，\(f(\mu)\) 表示它在循环型 \(\mu\vdash n\) 上的值，\(z_\mu\) 为该循环型的中心化子阶，\(p_\mu\in\Lambda_{\mathbb C}^n\) 为幂和对称函数的分拆乘积。定义
	\[
		\operatorname{ch}_n(f) =\sum_{\mu\vdash n}\frac{f(\mu)}{z_\mu}p_\mu.
	\]
	从复类函数空间 \(\operatorname{CF}(S_n)\) 到复对称函数的 \(n\) 次齐次分量 \(\Lambda_{\mathbb C}^n\) 的这个线性映射称为\term[Frobenius-tezheng-yingshe]{Frobenius 特征映射}[Frobenius characteristic map]。对有限维复 \(S_n\)-表示 \(U\)，也记 \(\operatorname{ch}(U)=\operatorname{ch}_n(\chi_U)\)。
\end{definition}
\symidx{\(\operatorname{ch}_n\)：Frobenius 特征映射}

这里的 \(z_\mu\) 与\cref{b5:lem:schur-cauchy}一致。事实上，与一个循环型为 \(\mu\) 的置换交换的置换，可以交换等长轮换，再在每个轮换内部独立转动。因此其中心化子阶为 \(\prod_r r^{m_r(\mu)}\cdot m_r(\mu)!\)，共轭类大小为 \(n!/z_\mu\)。

\begin{proposition}\label{b5:prop:frobenius-isometry}
	对任意非负整数 \(n\)，Frobenius 特征映射 \(\operatorname{ch}_n:\operatorname{CF}(S_n)\to\Lambda_{\mathbb C}^n\) 是复向量空间同构。对任意复类函数 \(f,g\in\operatorname{CF}(S_n)\)，以第一变量共轭线性的 Hall 内积配对它们的像，有
	\[
		\bigl\langle\operatorname{ch}_n(f),\operatorname{ch}_n(g)\bigr\rangle =\frac1{n!}\sum_{\sigma\in S_n}\overline{f(\sigma)}g(\sigma) =\sum_{\mu\vdash n}\frac{\overline{f(\mu)}g(\mu)}{z_\mu}.
	\]
	这里 \(f(\mu),g(\mu)\) 为循环型 \(\mu\) 上的值，\(z_\mu\) 为相应中心化子阶。并且平凡表示与符号表示的像分别为完全对称函数 \(h_n\)、初等对称函数 \(e_n\)。
\end{proposition}
\begin{proof}
	各循环型的示性函数构成类函数基，像是非零倍数的 \(p_\mu\)，所以映射为同构。Hall 内积对幂和基的对角公式立即给出所写等式；共轭类大小解释了两种求和形式。

	平凡表示的像为 \(\sum_{\mu\vdash n}p_\mu/z_\mu\)。把所有 \(n\) 汇总，其生成函数为
	\[
		\prod_{r\ge1}\left(\sum_{a\ge0}\frac{p_r^a t^{ra}}{r^a a!}\right) =\exp\left(\sum_{r\ge1}\frac{p_r}{r}t^r\right)=H(t).
	\]
	符号在循环型 \(\mu\) 上的值为 \(\prod_r(-1)^{(r-1)m_r(\mu)}\)。相应地把上式中 \(p_r\) 换成 \((-1)^{r-1}p_r\)，得到 \(E(t)\)，所以次数 \(n\) 项为 \(e_n\)。
\end{proof}

普通表示环 \(R_{\mathbb C}(S_n)\) 的乘法来自同一个群上的张量积。本节还需要另一种乘法，它会改变对称群的次数；这两种运算不能混用。

\subsection{诱导乘积与对称函数乘法}
\label{b5:subsec:0702:02}

\begin{definition}\label{b5:def:induction-product}
	设 \(R_{\mathbb C}(S_n)\) 为对称群 \(S_n\) 的复表示环，\(\mathcal R=\bigoplus_{n\ge0}R_{\mathbb C}(S_n)\)。给定非负整数 \(a,b\)、有限维复 \(S_a\)-表示 \(U\) 和 \(S_b\)-表示 \(W\)，记 \([U],[W]\) 为其在表示环中的类，\(U\boxtimes W\) 为 \(S_a\times S_b\) 在 \(U\otimes_{\mathbb C}W\) 上按各因子分别作用的表示。把 \(S_a\times S_b\) 按前 \(a\) 个与后 \(b\) 个字母的置换嵌入 \(S_{a+b}\)，定义
	\[
		[U]*[W] =\bigl[\operatorname{Ind}_{S_a\times S_b}^{S_{a+b}}(U\boxtimes W)\bigr],
	\]
	并对虚表示双线性延拓。这称为\term[youdao-chengji]{诱导乘积}[induction product]。次数零的平凡表示是单位。
\end{definition}
\symidx{\(\mathcal R,[U]*[W]\)：对称群的分次表示环及诱导乘积}

\begin{proposition}\label{b5:prop:frobenius-induction-product}
	设 \(a,b\) 为非负整数，\(f,g\) 分别为 \(S_a,S_b\) 的复类函数，\((f\boxtimes g)(\sigma,\tau)=f(\sigma)g(\tau)\)，\(S_a\times S_b\le S_{a+b}\) 按两个连续字母块嵌入。对 Frobenius 特征映射，有
	\[
		\operatorname{ch}_{a+b} \bigl(\operatorname{Ind}_{S_a\times S_b}^{S_{a+b}}(f\boxtimes g)\bigr) =\operatorname{ch}_a(f)\operatorname{ch}_b(g).
	\]
	因而诱导乘积使 \(\mathcal R=\bigoplus_{n\ge0}R_{\mathbb C}(S_n)\) 成为分次交换环，Frobenius 特征映射保持该乘法。
\end{proposition}
\begin{proof}
	固定循环型 \(\nu\vdash a+b\) 的置换 \(\sigma\)。陪集空间 \(S_{a+b}/(S_a\times S_b)\) 可以看作全部 \(a\) 元子集。诱导表示按这些陪集分为向量空间直和，\(\sigma\) 把每个陪集分量送到相应的另一分量；不固定的分量不贡献矩阵的对角块，因而不贡献迹。所以诱导特征标的求和只有在 \(\sigma\) 保持该子集时才有贡献。若一个数字属于被保持的子集，其所有 \(\sigma\) 迭代像也都属于该子集；因此这样的子集必为若干完整轮换之并。若它及其补集的循环型分别为 \(\alpha,\beta\)，则
	\[
		m_r(\alpha)+m_r(\beta)=m_r(\nu),\qquad |\alpha|=a,\quad|\beta|=b,
	\]
	而选择这些轮换的方法数为
	\[
		\prod_r\binom{m_r(\nu)}{m_r(\alpha)} =\frac{z_\nu}{z_\alpha z_\beta}.
	\]
	每种选择贡献 \(f(\alpha)g(\beta)\)。所以诱导类函数在 \(\nu\) 处的值除以 \(z_\nu\)，等于
	\[
		\sum_{\alpha\cup\beta=\nu} \frac{f(\alpha)g(\beta)}{z_\alpha z_\beta}.
	\]
	这正是右端幂和乘积中 \(p_\nu\) 的系数，证明了等式。类函数上的特征映射单射，而右端乘法结合、交换且有次数零的单位，因此同样性质在虚表示上成立。
\end{proof}

例如取 \(a=b=2\)，并令 \(U=W=1_{S_2}\)。诱导表示是 \(S_4\) 在全部二元子集上的置换表示。对 \(\sigma=(12)(34)\)，不变二元子集只能是完整轮换的支集 \(\{1,2\}\)、\(\{3,4\}\)，所以诱导特征标在 \(\sigma\) 处的值为二。这里
\[
	\nu=(2,2),\qquad \alpha=\beta=(2),\qquad z_\nu=2^2\cdot2!=8,\qquad z_\alpha=z_\beta=2,
\]
故选择数 \(z_\nu/(z_\alpha z_\beta)=2\) 正是在两个二轮换中选择一个放入第一个子集。另一方面，
\[
	\operatorname{ch}(1_{S_2})=h_2=\frac{p_1^2+p_2}{2},\qquad h_2^2=\frac{p_1^4+2p_1^2p_2+p_2^2}{4}.
\]
其中 \(p_{(2,2)}=p_2^2\) 的系数为 \(1/4=2/8\)，恰好是诱导特征标值除以 \(z_{(2,2)}\)。两个不变子集对应两种轮换分配，而特征映射的归一化把这个计数变成幂和乘积的系数。

\ProseParagraphBreak

设 \(\mu=(\mu_1,\ldots,\mu_r)\vdash n\)，Young 置换模为
\[
	M^\mu=\operatorname{Ind}_{S_{\mu_1}\times\cdots\times S_{\mu_r}}^{S_n}1.
\]
由刚证明的乘法公式及 \(\operatorname{ch}(1_{S_a})=h_a\)，有
\begin{equation}\label{b5:eq:young-permutation-characteristic}
	\operatorname{ch}(M^\mu)=h_\mu.
\end{equation}
同样的记法可用于任意非负整数组合：重排各部分只改变块的排列，诱导表示同构；零部分可以删去。

\begin{theorem}\label{b5:thm:frobenius-characteristic}
	设 \(n\) 为非负整数，\(\lambda\vdash n\)，\(S^\lambda\) 为对应的复 Specht 左 \(\mathbb C[S_n]\)-模，\(s_\lambda\) 为 Schur 函数，\(\operatorname{ch}_n\) 为 Frobenius 特征映射，记 \(\operatorname{ch}(S^\lambda)=\operatorname{ch}_n(\chi_{S^\lambda})\)。在有理系数对称函数代数中，记 Schur 函数的整系数线性组合所成的加法子群为
	\[
		\Lambda_{\mathbb Z}\coloneqq\bigoplus_\lambda\mathbb Zs_\lambda \subseteq\Lambda_{\mathbb Q},
	\]
	其中直和遍历全部分拆。则
	\begin{equation}\label{b5:eq:specht-schur}
		\operatorname{ch}(S^\lambda)=s_\lambda.
	\end{equation}
	这个加法子群 \(\Lambda_{\mathbb Z}\) 是 \(\Lambda_{\mathbb Q}\) 的分次子环，直和映射 \(\operatorname{ch}=\bigoplus_{n\ge0}\operatorname{ch}_n\) 给出带诱导乘积的 \(\mathcal R=\bigoplus_{n\ge0}R_{\mathbb C}(S_n)\) 到 \(\Lambda_{\mathbb Z}\) 的分次环同构，并将不可约表示的类对应到 Schur 基。
\end{theorem}
\begin{proof}
	取 \(N\ge n\)，将 \(\lambda\) 补零到 \(N\) 项。由 Jacobi--Trudi 公式和\eqref{b5:eq:young-permutation-characteristic}，\(s_\lambda\) 是下列虚特征标的像：
	\[
		\theta_\lambda =\sum_{\sigma\in S_N}\operatorname{sgn}(\sigma)\, \chi_{M^{\alpha(\sigma)}}, \qquad \alpha_i(\sigma)=\lambda_i-i+\sigma(i).
	\]
	含负部分的项规定为零；其他项将 \(\alpha(\sigma)\) 重排成分拆 \(\mu(\sigma)\)。

	对任意 \(k\)，有
	\[
		\sum_{i=1}^k\alpha_i(\sigma) =\sum_{i=1}^k\lambda_i+ \sum_{i=1}^k\sigma(i)-\frac{k(k+1)}2 \ge\sum_{i=1}^k\lambda_i.
	\]
	重排后前 \(k\) 项是原列中最大的 \(k\) 项，其和至少等于任取 \(k\) 项的和，尤其有
	\[
		\sum_{i=1}^k\mu_i(\sigma) \ge\sum_{i=1}^k\alpha_i(\sigma) \ge\sum_{i=1}^k\lambda_i.
	\]
	因此 \(\mu(\sigma)\) 支配 \(\lambda\)。若 \(\sigma\ne1\)，存在一个 \(k\) 使 \(\bigl\{\sigma(1),\ldots,\sigma(k)\bigr\}\ne\{1,\ldots,k\}\)，此时上式严格，故 \(\mu(\sigma)\ne\lambda\)。

	现在应用\cref{b5:prop:specht-dominance}：\(M^\mu\) 只含支配 \(\mu\) 的 Specht 类型，并且 \(S^\mu\) 的重数是一。因此
	\[
		\theta_\lambda=\chi_{S^\lambda} +\sum_{\eta\triangleright\lambda}c_\eta\chi_{S^\eta}, \qquad c_\eta\in\mathbb Z,
	\]
	其中 \(\triangleright\) 表示严格支配。另一方面，特征映射等距且 Schur 函数范数为一，故
	\[
		1=\langle s_\lambda,s_\lambda\rangle =\langle\theta_\lambda,\theta_\lambda\rangle =1+\sum_{\eta\triangleright\lambda}c_\eta^2.
	\]
	最后一式使用的是已经在\cref{b5:thm:specht-classification}建立的不可约分类与特征标正交性。于是全部 \(c_\eta=0\)，得到 \(\theta_\lambda=\chi_{S^\lambda}\)。不可约表示的类构成各 \(R_{\mathbb C}(S_n)\) 的 \(\mathbb Z\)-基，所以 \(\operatorname{ch}\) 给出 \(\mathcal R\) 到 \(\Lambda_{\mathbb Z}\) 的加法群同构。由\cref{b5:prop:frobenius-induction-product}，它还保持乘法和单位，故 \(\Lambda_{\mathbb Z}\) 对乘法封闭，所述映射为分次环同构。
\end{proof}

\begin{example}\label{b5:exm:frobenius-cancellation-22}
	在复数域上取 \(\lambda=(2,2)\)。两行 Jacobi--Trudi 行列式为 \(h_2^2-h_1h_3\)，对应的虚特征标是
	\[
		\theta_{(2,2)}=\chi_{M^{(2,2)}}-\chi_{M^{(3,1)}}.
	\]
	用前一章的 Young 子群不动向量计算，可以直接看见被消去的类型。令
	\[
		H_{22}=\langle(12),(34)\rangle,\qquad H_{31}=\langle(12),(23)\rangle.
	\]
	由\cref{b5:thm:maschke}及\cref{b5:thm:frobenius-reciprocity}，不可约模 \(S^\eta\) 在 \(\operatorname{Ind}_H^{S_4}1\) 中的重数等于 \(\dim(S^\eta)^H\)：前者给出完全可约性，后者把从诱导模到 \(S^\eta\) 的 Hom 空间识别为从平凡 \(H\)-模到 \(S^\eta\) 的 Hom 空间。\Cref{b5:prop:specht-dominance}又说明，\(M^{(2,2)}\) 只可能含 \((4),(3,1),(2,2)\) 三种类型，\(M^{(3,1)}\) 只可能含前两种。

	平凡模的不动空间均为一维。对 \(S^{(3,1)}\)，用 \(v_i\) 表示第二行只有 \(i\) 的 \((3,1)\)-行表类。第一列上下数字为 \(a,b\) 时，列交错元为 \(v_b-v_a\)；这些坐标差生成
	\[
		W=\left\{\sum_{i=1}^4a_i v_i:\sum_{i=1}^4a_i=0\right\}.
	\]
	这就恢复了\subsecref{b5:subsec:0606:01}中的 \(S^{(3,1)}=W\) 模型。\(H_{22}\) 不动向量的坐标为 \((a,a,-a,-a)\)，\(H_{31}\) 不动向量的坐标为 \((a,a,a,-3a)\)，因此两处重数也均为一。对 \(S^{(2,2)}\)，\cref{b5:prop:specht-22-char-two}给出的整数标准基 \(e_1,e_2\) 在复数域上满足
	\[
		\rho((12))=\rho((34))=\begin{pmatrix}1&0\\-1&-1\end{pmatrix},\qquad
		\rho((23))=\begin{pmatrix}0&1\\1&0\end{pmatrix}.
	\]
	向量 \(xe_1+ye_2\) 被第一个矩阵固定当且仅当 \(x+2y=0\)，所以 \((S^{(2,2)})^{H_{22}}=\mathbb C(2e_1-e_2)\)；再要求第二个矩阵固定，就有 \(x=y\)，因而 \((S^{(2,2)})^{H_{31}}=0\)。完全可约性于是给出
	\[
		\begin{aligned}
			M^{(2,2)}&\cong S^{(4)}\oplus S^{(3,1)}\oplus S^{(2,2)},\\
			M^{(3,1)}&\cong S^{(4)}\oplus S^{(3,1)}.
		\end{aligned}
	\]
	两项特征标相减时，平凡类型与标准类型分别相消，只留下 \(\chi_{S^{(2,2)}}\)。一般分拆的行列式有更多正负项；定理证明中的支配三角性先确定目标类型的重数为一，再由范数为一排除所有剩余的较高类型。
\end{example}

Young 图转置交换行求和与列交错求和；在特征映射下，它也交换完全对称函数与初等对称函数。

\begin{proposition}\label{b5:prop:omega-sign-conjugation}
	设 \(\Lambda_{\mathbb Q}\) 为有理系数对称函数代数，\(p_r,h_r,e_r\) 分别为幂和、完全与初等对称函数，\(s_\lambda\) 为 Schur 函数。赋值
	\[
		\omega(p_r)=(-1)^{r-1}p_r\qquad(r\ge1)
	\]
	唯一延拓为保次数的 \(\mathbb Q\)-代数自同构 \(\omega:\Lambda_{\mathbb Q}\to\Lambda_{\mathbb Q}\)。它是 Hall 配对的等距对合，且对非负整数 \(r\) 与任意分拆 \(\lambda\)，有
	\[
		\omega(h_r)=e_r,\qquad \omega(e_r)=h_r,\qquad \omega(s_\lambda)=s_{\lambda'},
	\]
	其中 \(\lambda'\) 为共轭分拆。将 \(\omega\) 复线性延拓到 \(\Lambda_{\mathbb C}\) 后，对任意非负整数 \(n\) 和有限维复 \(S_n\)-表示 \(U\)，有
	\[
		\operatorname{ch}(U\otimes_{\mathbb C}\operatorname{sgn}) =\omega\bigl(\operatorname{ch}(U)\bigr),
	\]
	其中张量积取对角左 \(S_n\) 作用，\(\operatorname{sgn}\) 是符号表示。
\end{proposition}
\symidx{\(\omega\)：交换完全与初等对称函数的对合}
\begin{proof}
	由\cref{b5:prop:symmetric-generators}，各 \(p_r\) 自由生成 \(\Lambda_{\mathbb Q}\)，所以给定赋值唯一延拓为代数同态。它不改变生成元的次数，且两次作用于每个 \(p_r\) 都得到 \(p_r\)，故 \(\omega^2=1\)，从而是保次数的自同构。对分拆 \(\mu\)，有
	\[
		\omega(p_\mu)=(-1)^{|\mu|-\ell(\mu)}p_\mu.
	\]
	Hall 配对的幂和基两两正交，而这个符号的平方为一，因此 \(\omega\) 保持 Hall 配对；复化后同样保持第一变量共轭线性的 Hall 内积。

	对\eqref{b5:eq:hp-generating}中的 \(H(t)\) 逐项施加 \(\omega\)，由生成函数关系得到
	\[
		\omega(H(t)) =\exp\left(\sum_{r\ge1}\frac{(-1)^{r-1}p_r}{r}t^r\right) =H(-t)^{-1}=E(t).
	\]
	比较 \(t^r\) 的系数得 \(\omega(h_r)=e_r\)，再用 \(\omega^2=1\) 得 \(\omega(e_r)=h_r\)。

	循环型 \(\mu\vdash n\) 的符号是 \((-1)^{n-\ell(\mu)}\)，所以 \(U\otimes\operatorname{sgn}\) 的特征标在该循环型上的值为 \((-1)^{n-\ell(\mu)}\chi_U(\mu)\)。逐项代入特征映射定义，恰好是在 \(\operatorname{ch}(U)\) 的每个 \(p_\mu\) 上施加 \(\omega\)，得到符号扭曲公式。最后，把\cref{b5:prop:specht-conjugate-sign}用于 \(k=\mathbb C\)，再用\eqref{b5:eq:specht-schur}，便有
	\[
		\omega(s_\lambda) =\omega\bigl(\operatorname{ch}(S^\lambda)\bigr) =\operatorname{ch}(S^\lambda\otimes\operatorname{sgn}) =\operatorname{ch}(S^{\lambda'})=s_{\lambda'}.
	\]
\end{proof}

例如 \((3,1)'=(2,1,1)\)。对 \(s_{(3,1)}=h_3h_1-h_4\) 施加 \(\omega\)，得到
\[
	s_{(2,1,1)}=e_3e_1-e_4.
\]
它与 \(S^{(3,1)}\otimes\operatorname{sgn}\cong S^{(2,1,1)}\) 表达同一个转置关系。

\subsection{特征标计算的组合化}
\label{b5:subsec:0702:03}

\begin{corollary}[Frobenius 特征标公式]\label{b5:cor:frobenius-coefficient}
	设 \(n\) 为非负整数，\(\lambda,\mu\vdash n\)，\(S^\lambda\) 为复 Specht 模，\(\chi_{S^\lambda}(\mu)\) 为其特征标在循环型 \(\mu\) 上的值，\(z_\mu\) 为相应中心化子阶。取整数 \(N\ge\ell(\lambda)\)，把 \(\lambda\) 补零至 \(N\) 项，令 \(x=(x_1,\ldots,x_N)\)、\(\delta_N=(N-1,\ldots,0)\)、\(\Delta_N(x)=\prod_{i<j}(x_i-x_j)\)。对这些变量中的 Schur 多项式和幂和多项式，有
	\[
		p_\mu(x)=\sum_{\eta\vdash n}\chi_{S^\eta}(\mu)s_\eta(x), \qquad \chi_{S^\lambda}(\mu) =[x^{\lambda+\delta_N}]\Delta_N(x)p_\mu(x).
	\]
	等价地，\(s_\lambda\) 中 \(p_\mu\) 的系数乘以 \(z_\mu\)，就是 \(\chi_{S^\lambda}(\mu)\)。
\end{corollary}
\begin{proof}
	由\eqref{b5:eq:specht-schur}，\(\langle p_\mu,s_\lambda\rangle=\chi_{S^\lambda}(\mu)\)。这些特征标值为实数：一个置换与其逆具有相同循环型，而有限群特征标满足 \(\chi(g^{-1})=\overline{\chi(g)}\)。所以交换内积的两个参数不会改变该值。在对称函数代数中按 Schur 标准正交基展开 \(p_\mu\)，再特化到 \(N\) 个变量，得第一式；长度超过 \(N\) 的分拆对应的 Schur 多项式为零。乘以 \(\Delta_N\) 后，其余各项为不同严格递减指数列的交错式。由于 \(N\ge\ell(\lambda)\)，目标项仍存在，且只有它在 \(x^{\lambda+\delta_N}\) 处贡献一，读取这个系数便得第二式。
\end{proof}

例如直接计算 \(S^{(2,1)}\) 在三轮换上的特征标。取 \(N=2\)，则 \(\delta_2=(1,0)\)、\((2,1)+\delta_2=(3,1)\)，并有
\[
	\Delta_2(x,y)p_3(x,y) =(x-y)(x^3+y^3)=x^4-x^3y+xy^3-y^4.
\]
读取 \(x^3y\) 的系数，得到 \(\chi_{S^{(2,1)}}((123))=-1\)。这里的负值直接来自 Vandermonde 因子中的 \(-y\) 项，无须先计算 \(s_{(2,1)}\) 的幂和展开。

\ProseParagraphBreak

例如
\[
	h_2=\frac{p_1^2+p_2}{2},\qquad h_3=\frac{p_1^3+3p_1p_2+2p_3}{6}, \qquad s_{(2,1)}=\frac{p_1^3-p_3}{3}.
\]
由于 \(z_{(1^3)}=6\)、\(z_{(2,1)}=2\)、\(z_{(3)}=3\)，立即读出 \(S_3\) 的二维特征标为 \((2,0,-1)\)。

\ProseParagraphBreak
对四次情形，
\[
	s_{(2,2)}=h_2^2-h_1h_3 =\frac{p_1^4}{12}+\frac{p_2^2}{4}-\frac{p_1p_3}{3}.
\]
按循环型 \((1^4),(2,1,1),(2,2),(3,1),(4)\) 排列，特征标值为 \((2,0,2,-1,0)\)。同一个行列式因而同时编码维数、换位上的迹以及较长轮换上的迹。

\ProseParagraphBreak
若实际计算 \(S^\lambda\) 在指定循环型上的特征标，可以先用 Jacobi--Trudi 求 \(s_\lambda\)，再用\eqref{b5:eq:newton-h}逐个把所需 \(h_r\) 化成幂和，最后读取一个系数。不需要预先求整张特征标表。Etingof 等人的叙述在\cite[Sections 5.14--5.15, pp. 114--118]{EtingofEtAl2011RepresentationTheory}中以 Young 置换模、三角性及 Cauchy 行列式建立相应公式；这里已把等距性与诱导乘法一并写出，以便继续处理张量分解。
